% Source: https://artofproblemsolving.com/community/c3h1881929_polynomial
% Stable identity: aops-topic-1881929-q01
Let the roots be $a$ and $b$. By Vieta's formulas $a+b=pq$ and $ab=p+q$. Both are positive, so $a$ and $b$ are positive integers. Now
\[
(a-1)(b-1)+(p-1)(q-1)=\bigl(ab-a-b+1\bigr)+\bigl(pq-p-q+1\bigr)=(p+q)-pq+1+pq-(p+q)+1=2 .
\]
Both products are nonnegative integers, so $(p-1)(q-1)\in\{0,1,2\}$.
\begin{itemize}
\item $(p-1)(q-1)=0$: say $p=1$. Then $(a-1)(b-1)=2$, so $\{a,b\}=\{2,3\}$, $q=a+b=5$, and indeed $p+q=6=ab$. This gives $(1,5)$ and $(5,1)$; the equation is $x^2-5x+6=0$ with roots $2,3$.
\item $(p-1)(q-1)=1$: $p=q=2$, and $(a-1)(b-1)=1$ gives $a=b=2$. The equation $x^2-4x+4=0$ has the double root $2$.
\item $(p-1)(q-1)=2$: $\{p,q\}=\{2,3\}$, and $(a-1)(b-1)=0$ with $a+b=6$ gives $\{a,b\}=\{1,5\}$. The equation $x^2-6x+5=0$ has roots $1,5$.
\end{itemize}
So the pairs are
\[
\boxed{(p,q)=(1,5),\ (5,1),\ (2,2),\ (2,3),\ (3,2)}
\]
(the pair $(2,2)$ gives a double root).
